\section{Diagnostic}
\label{sec:diagnostic}

Equation (12) of \citet{whitaker2025} freezes the hazard at the left endpoint of the remaining interval. The mean of the remaining incidence is that frozen hazard times the remaining time. \citet[Section 3.1]{golightly2015} use the same constant-hazard Gaussian step for remaining reaction counts. \citet[Section 3.2]{golightly2019} say that this step is often unreasonable, and their Section 6 says that accuracy falls as the inter-observation time grows. \citet[p.~25]{whitaker2025} call linearity of the latent process, over an interval no longer than the observation interval, a further assumption in the derivation of the conditioned hazard. Their Section 3.4 uses neither the word linear nor a second condition beyond equation (12) and the display after (18). We take the constant-hazard step to be the assumption the limitations paragraph is pointing at.

The same paragraph says the conditioned hazard is likely to propose trajectories inconsistent with the true conditioned process, when cumulative incidence does not evolve linearly. The law named on the page is the true conditioned process, and the sentence is about the trajectories the hazard proposes. The discretised transition in the numerator of Equation~\eqref{eq:weight} stays in the importance target, which for one propagate step is that transition times the observation mass function. This is a reading of the printed weight.

A poor \(h^{*}\) changes \(q\) and therefore the weights. The limitations paragraph gives no numerical spacing at which the 2025 guide fails for this epidemic model. \citet[Section 5.1]{golightly2019} report that a linear form fails to describe exponential decay, and that effective sample size falls as the interval length grows, for a death process. Those figures concern a death process, and the comparison below is for the epidemic model.

The observable is the effective sample size of the unnormalised incidence weights in one window,
\begin{equation}
\label{eq:ess}
\mathrm{ESS}(\tilde w_{1:N})
= \frac{\bigl(\sum_{i=1}^{N} \tilde w_i\bigr)^2}{\sum_{i=1}^{N} \tilde w_i^2},
\end{equation}
as printed by \citet[Section 5.2]{golightly2019}. Zero weights are included. If every weight in a window is zero, Equation~\eqref{eq:ess} is undefined, and the value entered in the table is 0. A draw has weight zero when the correction in Equation~\eqref{eq:hstar} is undefined, when a component of the hazard passed to the Poisson draw is negative, or when the updated counts would be negative. \citet{whitaker2025} print no truncation of \(h^{*}\). \citet[Section 3.2]{golightly2019} truncate their equation (5) to keep it positive. A negative hazard in the 2025 display is left untruncated and receives weight zero.

The comparison, specified before any weight was computed, is between the median of this effective sample size across replicate data sets at the coarsest spacing and the same median at the published unit spacing. A myopic proposal, with \(h^{*}\) replaced by \(h\), is computed on the same paths as a comparator. The diagnostic is that comparison across the three spacings. It supplies no threshold that would divide adequate spacing from inadequate spacing, and it leaves the filtering distribution unexamined. Section~\ref{sec:residual} states the proposal used in the second comparison.
